\documentclass{mathos}

\begin{document}
\renewcommand\bibname{Literatura}

\sloppy

\begin{titlepage}
\begin{center}
%\vspace*{-10ex}
    \includegraphics[keepaspectratio=true,scale=0.4]{MathosLogo.png} \\ \vspace{1mm}
       \textsc{Sveučilište Josipa Jurja Strossmayera u Osijeku} \\
    \vspace{2mm}
    \textsc{Fakultet primijenjene matematike i informatike}
\end{center}

\begin{center}

    \vspace{4mm}
     {Studij} %% Ispod su navedeni studiji i smjerovi

%% Diplomski sveučilišni nastavnički studij matematike i informatike
%% Sveučilišni diplomski studij matematike \\ modul: financijska matematika i statistika
%% Sveučilišni diplomski studij matematike \\ modul: matematika i računarstvo

\end{center}

\vspace{25mm}
\begin{center}
    % naslov
    {\LARGE{\bf  Naslov rada}}

    \vspace{15mm}
    %
    {\large{\sc Diplomski rad}}


\end{center}

\vspace{50mm}

\begin{minipage}[t]{0.47\textwidth}
	{Mentor:} \normalsize\vspace{3mm} %% Ako je potrebno dodati komentora odkomentirajte redove 45, 46, 47. Pogledajte donje primjere za navođenje imena mentora/komentora:
	{\bf\\ \large{Ime i prezime mentora}}
%% \\[2ex]
%% {Komentor:} \normalsize\vspace{3mm}
	%% {\bf\\ \large{Ime i prezime komentora}}


%% Imena mentora/komentora:

%% prof.\,dr.\,sc. Ime i prezime
%% izv.\,prof.\,dr.\,sc. Ime i prezime
%% doc.\,dr.\,sc. Ime i prezime
%% dr.\,sc. Ime i prezime
\end{minipage}
\hfill
\begin{minipage}[t]{0.47\textwidth}\raggedleft
	{{Student:}{\normalsize\vspace{3mm} \bf\\ \large{Ime i prezime studenta}}}
\end{minipage}

\vspace*{\fill}
\centering{\normalsize{Osijek, 20xx}}
\end{titlepage}


\tableofcontents
\thispagestyle{empty}


\chapter{Uvod}\label{Uvod}
\pagenumbering{arabic}

\noindent
Diplomski rad sadrži sljedeće cjeline i to navedenim redoslijedom:
\bigskip

{\it Sadržaj}: sastoji se od naslova poglavlja i točaka s rednim brojevima odgovarajućih početnih
stranica.

{\it Uvod}: treba sažeto opisati diplomski zadatak te navesti sadržaj pojedinih
poglavlja diplomskog rada.

{\it Glavni dio rada}: treba ga podijeliti u poglavlja i točke koje
čine određenu logičku cjelinu. Moguće je dodati i odgovarajuće potpoglavlje.
Neki savjeti za pisanje ovog dijela rada dani su i u Poglavljima~\ref{poglavlje_1} i \ref{poglavlje_2} ovih uputa.

{\it Literatura}: korištenu literaturu treba navesti abecednim redom prema prezimenu prvog autora.
Primjeri nekih referenci nalaze se u poglavlju {\it Literatura}. Korištene reference potrebno je
citirati na odgovarajućim mjestima u radu.

{\it Sažetak}: ukratko treba opisati sadržaj diplomskog rada.

{\it Ključne riječi}: navesti listu ključnih riječi vezanih uz sadržaj rada.

{\it Engleski prijevod}: title, summary i keywords odnose se na englesku verziju naslova rada,
sažetak i ključne riječi.

{\it Životopis}: piše se u prvom licu jednine i treba sadržavati biografske podatke,
podatke o sudjelovanju na natjecanjima te podatke o osvojenim nagradama.

\bigskip
\noindent
Obzirom na naslovnu stranicu, pažljivo pogledajte zakomentirane dijelove LaTeX koda, a odnose se na smjer studija, pisanje imena mentora i eventualno dodavanje komentora.
Svako novo poglavlje ({\verb $\$chapter}) počinje na novom listu. Ako vam je potreban neki dodatni paket, jednostavno ga dodajte na početku koda.

\bigskip
\noindent
Predlaže se dvostrano štampanje diplomskog rada pa nije potrebno brisati prazne stranice generirane ovim kodom.


\chapter{Upute - prvi dio}\label{poglavlje_1}

\section{Matematička okruženja}

Definicije, teoreme i ostale rezultate s mogućim dokazima pišemo koristeći odgovarajuća okruženja. Primjerice,
\begin{defin}
Ovo je prva definicija.
\end{defin}
Ako teorem ili bilo koji drugi rezultat nije vaš originalni rezultat, onda ga treba citirati. Dakle, napisati odakle je taj rezultat preuzet. Primjerice,
\begin{teorem}[{vidjeti \cite[Theorem 1]{Acz}}]  \label{teorem1}
Neka je $n$ prirodan broj. \dots
\end{teorem}

\begin{proof}
Ovdje pošemo dokaz Teorema~\ref{teorem1}.
\end{proof}
Uočite kako je simbol $\Box$ po automatizmu dodan na kraju dokaza i ne treba ga ``umjetno'' dodavati.

\begin{primjer}\label{prim1}
Treba nastojati samostalno konstruirati primjere. Ako su neki od njih također preuzeti iz literature, opet ih treba citirati poput Teorema~\ref{teorem1}.
\end{primjer}
\noindent
{\it Rješenje.} Ovdje dolazi rješenje vašeg Primjera~\ref{prim1} (ako je predviđeno u tome obliku).

\subsection{Tablice i slike}

Slike uglavnom stvaramo u {\it .eps} formatu. Može se koristiti i drugi format slike, kao primjerice  Mathos logo na naslovnoj stranici, u formatu {\it .png}. Slike uvodimo na sljedeći način:

\begin{figure}[h!]
\centering
\includegraphics[scale=0.9]{fig1.eps}
\caption{Ovo je naša slika.}
\label{fig1}
\end{figure}

Ako je Slika~\ref{fig1} preuzeta iz neke literature, također ju treba citirati.

\newpage
Tablice se uvode na standardni način, a oblik ovisi o njihovoj strukturi. Primjerice,

\begin{table}[hbtp]
\centering
\begin{tabular}{|c|c|c|}\hline
Red & Stupac& Ukupno\\\hline
2 & 3 & 5\\\hline
3 & 9 & 12\\\hline
\end{tabular}
\caption{Ovo je prva tablica.}
\label{table1}
\end{table}

\noindent
Na rezultate, slike, tablice i ostale strukture u tekstu se poziva na sljedeći način, ovisno o kontekstu rečenice: ``Prema dokazu Teorema~\ref{teorem1} slijedi \dots''; ``na
Slici~\ref{fig1} prikazano je \dots''; ``iz Tablice~\ref{table1} može se zaključiti \dots''; ``Više o \dots može se vidjeti u~\cite{JunBae}''.



\chapter{Upute - drugi dio}\label{poglavlje_2}

\section{Prijelomi formula}
Kod korištenja okruženja za pisanje formula treba voditi računa o njihovom prijelomu i numeriranju. Evo nekoliko primjera:
\begin{align*}
    &u+\arctan x +x^3+4x^2+\sin{x}  \\
    &\qquad =\frac{1}{1 + x^2} \sin{2x}+e^x.
\end{align*}

\begin{eqnarray}\label{eq_1}
\begin{aligned}
    &u+\arctan x +x^3+4x^2+\sin{x}  \\
    &\qquad =\frac{1}{1 + x^2} \sin{2x}+e^x.
\end{aligned}
\end{eqnarray}

\begin{eqnarray}
\begin{aligned}
    u=&\arctan x +x^3+4x^2+\sin{x}  \\
    &+\frac{1}{1 + x^2} \sin{2x}+e^x.
\end{aligned}
\end{eqnarray}

\begin{align*}
    u=&\arctan x +x^3+4x^2+\sin{x}  \\
    v=&\frac{1}{1 + x^2} \sin{2x}+e^x.\label{eef}
\end{align*}

\begin{align}
    u=&\arctan x +x^3+4x^2+\sin{x}  \\
    v=&\frac{1}{1 + x^2} \sin{2x}+e^x.\label{eef}
\end{align}

\begin{align*}
    u=&\arctan x +x^3+4x^2+\sin{x}  \\
    &+\frac{1}{1 + x^2} \sin{2x}+e^x\\
    =&\arctan x +x^3+4x^2+\sin{x}.
\end{align*}
Nije potrebno numerirati svaki izraz već samo one na koje ćete se u tekstu pozivati. Pozivamo se na način: ``Prema formuli \eqref{eq_1} \dots''

\noindent
Uočite kako formule i ostali matematički izrazi nikada ne završavaju s $+,-,\cdot,=$ već se znakovi prenose u novi red, kao u gornjim primjerima.

\noindent
Kod pisanja rada treba paziti da tekst, formule, tablice, slike, itd., ne izlaze iz margina rada.

\section{Programerska okruženja}

U ovom radu možete pisati pseudokod korištenjem \verb+algpseudocode+ paketa.

\begin{algorithm}
\caption{primjer algoritma}
\begin{algorithmic}[1]
\Procedure{Algoritam}{$n$}
\State $i \gets n$
\If{$i\geq 5$}
    \State $i \gets i-1$
\Else
    \If{$i\leq 3$}
        \State $i \gets i+2$
    \EndIf
\EndIf
\EndProcedure
\end{algorithmic}
\end{algorithm}

\noindent
Postoji mogućnost izlistavanja programskog koda koristeći \verb+listing+ paket.

\lstinputlisting[language=C++]{code1.cpp}


%% Literatura

\begin{thebibliography}{30}
\addcontentsline{toc}{chapter}{Literatura}
\bibitem{Acz}
{\sc J.~Acz\'{e}l}, {\em Lectures of Functional Equations and Their Applications}, Academic Press, New York, 1966. ({\it ovo je korištena knjiga}).

\bibitem{JunBae}
{\sc S.~Jung, J.~Bae}, {\em  Some functional equations originating from number theory}, Proc. Indian Acad. Sci. (Math. Sci.) {\bf{113}}(2003), 91--98.
({\it ovo je korišteni znanstveni članak}).

\bibitem{ } {\em The On-Line Encyclopedia of Integer Sequences}, dostupno na \newline {\verb"\https://oeis.org"}. ({\it ovo je korišteni izvor s interneta}).
\end{thebibliography}




\chapter*{Sažetak}
\markboth{}{SAŽETAK}
\addcontentsline{toc}{chapter}{Sažetak}

Ovdje pišemo sažetak diplomskog rada.

\section*{Ključne riječi}

Ovdje malim slovima (osim ako se ne radi o imenu) redom pišemo ključne riječi.




\chapter*{Naslov rada na engleskom jeziku}
\markboth{}{SUMMARY}

\section*{Summary}
\addcontentsline{toc}{chapter}{Summary}

Ovdje pišemo sažetak na engleskom jeziku.

\section*{Keywords}

Ovdje malim slovima (osim ako se ne radi o imenu) redom pišemo ključne riječi na engleskom jeziku.




\chapter*{Životopis}
\addcontentsline{toc}{chapter}{Životopis}

Ovdje pišemo životopis u prvom licu jednine.

\end{document}
